\section*{Hoofdstuk \ref{ch:g1:measure} --- Meten en tijd}

\begin{solution}{exo:g1:measure:1}
Het antwoord hangt af van de potloden; met één uiteinde gelijk komt het
langste potlood het verst en het kortste het minst ver.
\end{solution}

\begin{solution}{exo:g1:measure:2}
De antwoorden verschillen; je leest elke lengte af door de $0$ van de
lat bij het begin van het voorwerp te leggen en aan het einde de hele
centimeter af te lezen.
\end{solution}

\begin{solution}{exo:g1:measure:3}
Na maandag: dinsdag. Voor zondag: zaterdag. Twee dagen na vrijdag:
zondag.
\end{solution}

\begin{solution}{exo:g1:measure:4}
Gisteren: woensdag. Morgen: vrijdag. Over drie dagen: zondag.
\end{solution}

\begin{solution}{exo:g1:measure:5}
Om $6$ uur wijst de korte wijzer naar de $6$ en de lange wijzer recht
naar boven, naar de $12$.
\end{solution}

\begin{solution}{exo:g1:measure:6}
Korte wijzer op $5$: $5$ uur. Korte wijzer op $11$: $11$ uur.
\end{solution}

\begin{solution}{exo:g1:measure:7}
Voor $8$ bijvoorbeeld $5 + 2 + 1$ of $2 + 2 + 2 + 2$ (er zijn nog meer
goede antwoorden).
\end{solution}

\begin{solution}{exo:g1:measure:8}
$5 + 2 + 2 + 2 = 11$. Leo heeft $11$.
\end{solution}

\begin{solution}{exo:g1:measure:9}
Van $3$ naar $5$ is $2$: Zoë krijgt $2$ terug.
\end{solution}
